\section{Latihan}
\label{sec:latihan-bab13}

\begin{enumerate}
    \item \textbf{Konseptual.} Sebutkan tiga alasan mengapa program
          sebaiknya dipecah menjadi prosedur. Jelaskan pula mengapa
          pemakaian prosedur tidak selalu membuat program menjadi
          lebih cepat, dan kapan prosedur sebaiknya dihindari.

    \item \textbf{Konseptual.} Jelaskan apa yang dikerjakan
          instruksi \asm{CALL} dan \asm{RET}. Mengapa \asm{JMP}
          tidak dapat dipakai sebagai pengganti \asm{CALL}?

    \item \textbf{Analisis stack.} Diketahui \reg{SS} $=$ \hex{2000}
          dan \reg{SP} $=$ \hex{FFFC}. Tentukan alamat fisik sel
          teratas stack. Kemudian tentukan nilai \reg{SP} dan isi
          sel tersebut setelah urutan instruksi berikut dijalankan:

\begin{lstlisting}[language=MASM, numbers=none]
    MOV AX, 1234h
    PUSH AX
    PUSH BX
    POP  CX
\end{lstlisting}

          \begin{enumerate}
              \item Berapa byte data yang masih tertinggal di stack
                    setelah instruksi terakhir?
              \item Apakah nilai \reg{BX} masih tersimpan di memori
                    stack? Jelaskan.
          \end{enumerate}

    \item \textbf{Analisis \eng{stack frame}.} Perhatikan potongan
          prosedur berikut.

\begin{lstlisting}[language=MASM, numbers=none]
Jumlah PROC
    PUSH BP
    MOV  BP, SP
    MOV  AX, [BP+4]
    ADD  AX, [BP+6]
    POP  BP
    RET  4
Jumlah ENDP
\end{lstlisting}

          \begin{enumerate}
              \item Tuliskan urutan instruksi pemanggil yang tepat
                    agar prosedur ini menghasilkan nilai benar.
              \item Nilai apa yang berada pada \code{[BP+8]} saat
                    prosedur berjalan?
              \item Mengapa pemanggil tidak boleh menambahkan
                    \lstasm|ADD SP, 4| setelah \asm{CALL}?
          \end{enumerate}

    \item \textbf{Analisis kesalahan.} Sebuah prosedur memakai
          register \reg{SI} dan \reg{DI} untuk mengolah data, tetapi
          hanya menyimpan \reg{BP} pada \eng{prologue}-nya. Setelah
          prosedur dipanggil, program pemanggil menghasilkan keluaran
          yang salah. Jelaskan penyebabnya dan tuliskan perbaikan
          yang diperlukan, dengan mengacu pada
          Tabel~\ref{tab:register-simpan}.

    \item \textbf{Penulisan program.} Tulislah prosedur
          \code{HitungRataRata} yang menerima dua parameter 16-bit
          melalui stack dan mengembalikan hasil bagi keduanya
          (dibulatkan ke bawah) pada register \reg{AX}. Sertakan
          \eng{prologue}, \eng{epilogue}, dan potongan program
          pemanggilnya.

    \item \textbf{Penulisan program.} Tulislah prosedur rekursif
          \code{Pangkat} yang menghitung $a^b$ untuk $b \ge 0$,
          dengan parameter $a$ dan $b$ dikirim melalui stack serta
          hasil dikembalikan pada \reg{AX}. Tentukan pula kasus
          basisnya dan jelaskan mengapa kasus basis itu diperlukan.

    \item \textbf{Perbandingan.} Bandingkan pemakaian stack pada
          prosedur faktorial rekursif dan faktorial iteratif untuk
          $n = \dec{10}$. Hitung perkiraan jumlah byte stack yang
          dipakai masing-masing versi, lalu simpulkan versi mana yang
          lebih tepat dipakai dan pada keadaan apa rekursi tetap
          bermanfaat.
\end{enumerate}

\begin{refleksibox}
Jawab pertanyaan berikut pada jurnal praktikum Anda.
\begin{enumerate}
    \item Bagian mana dari bab ini yang paling sulit dipahami:
          arah tumbuh stack, susunan \eng{stack frame}, atau konvensi
          pemanggilan? Apa yang Anda lakukan untuk mengatasinya?
    \item Setelah memeriksa isi stack pada \eng{debugger}, apakah
          gambaran Anda tentang ``pemanggilan prosedur'' berubah?
          Jelaskan bagian mana yang paling mengejutkan.
    \item Konsep apa yang masih perlu Anda latih sebelum masuk ke
          Bab~\ref{chap:interupsi}, yang juga memakai stack untuk
          menyimpan \reg{CS}, \reg{IP}, dan \reg{FLAGS}?
\end{enumerate}
\end{refleksibox}
